\section{Layers}\label{sec:layers}

The later sections use Theorem~\ref{thm:forest} through Corollary~\ref{cor:cycles} and through
one lemma.

\begin{idea}{the content of the lemma}
If the outside part is a forest for more than half of the even subgraphs, by a fixed margin, then
every distance layer around the marked set is small. In the other case a window of consecutive layers is a large set
with a small boundary, and the forest estimate applies inside this set. We use the lemma for
paths in Section~\ref{sec:poles} and for large girth in Section~\ref{sec:half}.
\end{idea}

Let $G$ be a multigraph and $P$ a set of vertices. A \emph{level function} is a map
$\lambda\colon V(G)\to\mathbb{Z}$ such that $|\lambda(u)-\lambda(w)|\le1$ for every edge $uw$ that
has an end outside $P$. The distance from a fixed vertex is a level function. For
$\delta\in(0,\half)$ let $R_\delta\ge2^{100}$ be an integer such that every $r\ge R_\delta$ is in
the range of Theorem~\ref{thm:forest} and satisfies $20/\lgk3r\le\delta$. For an integer $k\ge1$
put
\[
  \Theta_\delta(k)=\max\bigl\{2R_\delta,\;2^{100}(k+1)^{2}\bigr\}.
\]

\begin{lemma}[Layers]\label{lem:width}
Let $G$ be a connected subcubic multigraph, and let $P$ be a nonempty set of vertices such that
$G[P]$ is connected and every vertex outside $P$ has degree three. Let $\lambda$ be a level
function such that at most $k$ vertices of $P$ have any given level. Let $\delta\in(0,\half)$, and
suppose that $q(G,P)>\half+\delta$. Then, for every $j$, fewer than $\Theta_\delta(k)$ vertices
outside $P$ have level $j$.
\end{lemma}

\begin{proof}
Put $\Phi(m)=m/\log^{2}m$ and $\Theta=\Theta_\delta(k)$.

\step{Step 1: no large set outside $P$ has a small boundary.} We claim that no set
$X\subseteq V(G)\sm P$ has $|X|=m\ge\Theta$ and $|\delta X|\le\Phi(m)/4$. Suppose that $X$ is such a
set. Let $C_0$ be the component of $G-X$ that contains $P$, and put $X^{+}=V(G)\sm V(C_0)$. Then
$X\subseteq X^{+}$, the set $X^{+}$ avoids $P$, and $G-X^{+}=C_0$ is connected. Different components
of $G-X$ have no edge between them, so $\delta X^{+}\subseteq\delta X$. Every vertex of $X^{+}$ has
degree three. Put $d^{+}=|\delta X^{+}|\ge1$.

We form a multigraph $\Gamma$ from $G[X^{+}]$. For each edge of $\delta X^{+}$ add a new vertex
and join it to the end of that edge in $X^{+}$. Join the $d^{+}$ new vertices by a cycle $W$
through all of them, and let $P^{+}$ be the set of the new vertices. Every component of
$G[X^{+}]$ has an edge to $C_0$, so $\Gamma$ is connected. It is subcubic, $\Gamma[P^{+}]$ is
connected, the vertices outside $P^{+}$ have degree three, and $\Gamma-P^{+}=G[X^{+}]$. The
multigraph $\Gamma$ has $d^{+}$ more vertices and $2d^{+}$ more edges than $G[X^{+}]$, and
$3|X^{+}|=2|E(G[X^{+}])|+d^{+}$. So its rank is
\[
  \rho=\frac{|X^{+}|+d^{+}}2+1\;\ge\;\frac m2\;\ge\;R_\delta .
\]
The function $\Phi$ is increasing on $[8,\infty)$, and $\Phi(m)/4\le\Phi(m/2)$. So
$d^{+}\le\Phi(m)/4\le\Phi(m/2)\le\Phi(\rho)$, and Theorem~\ref{thm:forest} with the choice of
$R_\delta$ gives $q(\Gamma,P^{+})\le\half+\delta$.

The set $\xi\cap E(G[X^{+}])$ has the same distribution for $\xi$ uniform on $\Z(G)$ and for
$\xi$ uniform on $\Z(\Gamma)$. Indeed, contract the connected set $V(C_0)$ of $G$ to one vertex,
and delete the loops that arise. Call the result $G'$. An even subgraph of $G$ restricts to an
even subgraph of $G'$. Every even subgraph of $G'$ arises in this way: the missing edges inside
the contracted set can be chosen in a spanning tree of that set. So the restriction is a
surjective linear map $\Z(G)\to\Z(G')$, and it sends the uniform distribution to the uniform
distribution. The contraction of $W$ in $\Gamma$ gives the same multigraph $G'$, with the new
edges in the place of the edges of $\delta X^{+}$. Hence
\[
  \PP\bigl(\xi\cap E(G[X^{+}])\text{ is a forest}\bigr)\;=\;q(\Gamma,P^{+})\;\le\;\frac12+\delta .
\]
But $G[X^{+}]$ is a subgraph of $G-P$, so this probability is at least $q(G,P)>\half+\delta$. This
contradiction proves the claim.

\step{Step 2: a window of levels.} Let $m_j$ be the number of vertices outside $P$ of level $j$.
For $i\le i'$ let $X(i,i')$ be the set of vertices outside $P$ with a level in $[i,i']$. An edge
that leaves $X(i,i')$ ends outside $P$ at level $i-1$ or $i'+1$, or it ends in $P$ at a level in
$[i-1,i'+1]$. So
\begin{equation}\label{eq:window}
  |\delta X(i,i')|\;\le\;3m_{i-1}+3m_{i'+1}+3k(i'-i+3).
\end{equation}
Suppose that $m_j\ge\Theta$ for some $j$. Start with the window $i=i'=j$. As long as a level next
to the window has more than $m/(48\log^{2}m)$ vertices outside $P$, where $m=|X(i,i')|$ is the
present size, add that level to the window. The graph is finite, so this stops. Let $\mu$ be the
final size and $\gamma=\log\mu\ge100$. Each step multiplies the size by more than
$1+1/(48\gamma^{2})$. So the number of steps is at most
\[
  \frac{\ln\mu}{\ln\bigl(1+1/(48\gamma^{2})\bigr)}\;\le\;96\gamma^{2}\ln\mu\;<\;67\gamma^{3}.
\]
At the end both levels next to the window have at most $\mu/(48\gamma^{2})$ vertices outside $P$.
By~\eqref{eq:window},
\[
  |\delta X(i,i')|\;\le\;\frac{\mu}{8\gamma^{2}}+3k\bigl(67\gamma^{3}+3\bigr)
  \;\le\;\frac{\mu}{4\gamma^{2}}=\frac{\Phi(\mu)}4 .
\]
The second inequality is equivalent to $24k\gamma^{2}(67\gamma^{3}+3)\le2^{\gamma}$. This holds because
$k\le2^{\gamma/2-50}$, by $\mu\ge\Theta\ge2^{100}(k+1)^{2}$, and because
$24\gamma^{2}(67\gamma^{3}+3)\le2000\gamma^{5}\le2^{\gamma/2+50}$ for $\gamma\ge100$. So the final
window contradicts Step~1.
\end{proof}

